% Nouveau EGA V (cote 25): archive cote 25, site page 37 (PDF page 38).
% typescript. Draft compared: none.
% AI-transcribed by Claude Opus 5.5 (claude-opus-5-5, adaptive thinking, effort high), Message Batches API, batch ega25b (msgbatch_01A7n5YuXubUPNK1zsepi8RG), 2026-11-13; not yet checked by eye.
\scanpage{25/37}
\mspage{\hand{\unsure{39|33}}}

8.\over{3}{4}. Nous étudions maintenant les points  y de $\underline{Y}$ qui sont non lisses
sur $Y_{\hand{\zeta}}$ rel à $k(\hand{\zeta})$. Nous nous  bornons au cas où  $f:X\to P$  est \uline{non ramifié}
(pratiquement, ce sera une immersion), et où $X\to S$ est lisse.  Nous ne
supposons pas nécessairement S le spectre d'un corps.  \xout{Nous nous bornons}
\xout{à la considération des points de $X_P$ qui sont dans \unsure{l'ouvert U de}} Comme
f  est non ramifié, l'homomorphisme canonique   $f^*(\underline{\Omega}^1_{P/S})\to\underline{\Omega}^1_{X/S}$
est surjectif, et son noyau est un Module loc libre sur X, que nous
noterons  $\hand{\check{\nu}}_{X/P}$ ; lorsque f est une immersion, ce n'est autre que le
Module \uline{conormal}  $\underline{J}/\underline{J}^2$  défini par l'Idéal $\underline{J}$ de X dans P , et nous l'appe-%
lerons en tous cas Module conormal:
\[ \hand{(a)} \qquad 0\to\hand{\check{\underline{\nu}}}_{X/P}\to f^*(\underline{\Omega}^1_{P/S})\to\underline{\Omega}^1_{X/S}\to 0 \ . \]
\margin{$\|\|$ (vertical pencil strokes)} Notons que nous avons aussi sur  P  une suite exacte canonique
(qui pouvait figurer à titre d'exemple au par.16 p.ex.)
\[ \hand{(b)} \qquad 0\to\underline{\Omega}^1_{P/S}(1)\to\underline{E}_P\to\underline{O}_P(1)\to 0 \]
(i.e. $\underline{\Omega}^1_{P/S}$ est can isomorphe au noyau de l'homomorphisme canonique
$\underline{E}_P(-1)\to\underline{O}_P$  déduit de $\underline{E}_P\to\underline{O}_P(1)$ ), d'où en appliquant $f^*$ :
\[ \hand{(b')} \qquad 0 - f^*(\underline{\Omega}^1_{P/S})(1) - \underline{E}_X - \underline{O}_X(1)\to 0 \ \hand{,} \]
qui donne une description explicite de  $f^*(\underline{\Omega}^1_{P/S})$  \xout{\unsure{de X}} sur X, et
permet  donc d'identifier $\hand{\check{\underline{\nu}}}_{X/P}(1)$  à un sous-Module loc facteur direct
de $\underline{E}_X$ , ou encore le \xout{Module conormal} dual  $\hand{\underline{\nu}}_{X/P}(-1)$  à un Module
quotient de $\hand{\check{}}\underline{E}_X$ .  Par suite,   $P(\hand{\underline{\nu}}_{X/P}\deleted{\unsure{$*$}}(-1)\ )$  =  $P(\hand{\underline{\nu}}_{X/P})$  se plonge
canoniquement dans  $P(\hand{\check{}}\underline{E}_X) = X\times_S\hand{\check{}}P = X_{\hand{\check{}}P}$ comme un sous-fibré projectif sur
X, donc comme sous-préschéma fermé. Ce dernier est nécessairement contenu
dans $\underline{Y}$  (du fait que $\underline{\Omega}_{X/P}(1)$  est contenu dans le noyau de  $\underline{E}_X\to\underline{O}_X(1)$)
\deleted{L'ensemble sous-jacent  à ce sous-préschéma n'est autre que l'ensemble des}
\deleted{\xout{xxxxP} $y\in\overline{Y}$ \xout{\illegible{1}} \uline{qui sont} \inserted{\deleted{\unsure{des zéros x singuliers pour}}} \deleted{"singuliers sur $Y_{\hand{\zeta}}$ " au sens suivant:}}
\deleted{\margin{\unsure{des zéros singuliers}} \deleted{un point est dit "non singulier" sur $Y_{\hand{\zeta}}$ si $\hand{\phi_\zeta}$ est \unsure{régulier}} sur $X_{k(\hand{\zeta})}$}
\deleted{en y et \xout{\illegible{1}} $Y_{\hand{\zeta}}$ est lisse sur $k(\hand{\zeta})$ en y.}

\mspage{\hand{37}}
